\section{Experimental Results}
\label{sec:results}

Present your experimental setup and results.

\subsection{Experimental Setup}
\label{subsec:setup}

Describe the experimental environment, datasets, parameters, and evaluation metrics.

\begin{table}[h]
\centering
\caption{Experimental Parameters}
\label{tab:parameters}
\begin{tabular}{lcc}
\toprule
Parameter & Value & Description \\
\midrule
Parameter 1 & 100 & Description of parameter \\
Parameter 2 & 0.01 & Description of parameter \\
Parameter 3 & 50 & Description of parameter \\
\bottomrule
\end{tabular}
\end{table}

\subsection{Performance Evaluation}
\label{subsec:performance}

Present and analyze your results. Use tables and figures effectively.

\begin{figure}[h]
\centering
% Replace with actual figure: \includegraphics[width=0.5\columnwidth]{figure.pdf}
\framebox[0.5\columnwidth]{\rule{0pt}{3cm}Figure Placeholder}
\caption{Performance comparison of different methods}
\label{fig:performance}
\end{figure}

\subsection{Comparison with State-of-the-Art}
\label{subsec:comparison}

Compare your approach with existing methods. Provide quantitative analysis.

\begin{table}[h]
\centering
\caption{Performance Comparison}
\label{tab:comparison}
\begin{tabular}{lccc}
\toprule
Method & Metric 1 & Metric 2 & Metric 3 \\
\midrule
Method A & 85.3 & 12.4 & 0.82 \\
Method B & 87.1 & 11.8 & 0.85 \\
Method C & 89.2 & 10.5 & 0.88 \\
\textbf{Ours} & \textbf{92.4} & \textbf{9.2} & \textbf{0.91} \\
\bottomrule
\end{tabular}
\end{table}

\subsection{Ablation Study}
\label{subsec:ablation}

Present ablation studies to validate the contribution of different components.

\begin{table}[h]
\centering
\caption{Ablation Study}
\label{tab:ablation}
\begin{tabular}{lcc}
\toprule
Configuration & Accuracy & Speed (ms) \\
\midrule
Baseline & 85.3 & 12.4 \\
+ Component A & 87.6 & 13.1 \\
+ Component B & 89.2 & 13.8 \\
Full Model & 92.4 & 14.2 \\
\bottomrule
\end{tabular}
\end{table}
